\section{Experimental Setup}
\label{sec:setup}

\subsection{Dataset}
\label{sec:setup:dataset}

In addition to the Buy transactions and risk-band fields described in Sections~\ref{sec:intro} and~\ref{sec:method}, FAR-Trans~\cite{sanzcruzado2024fartransinvestmentdatasetfinancial} provides 159{,}136 Sell transactions and daily closing prices spanning January 2018 to November 2022. Of the 29{,}090 retail customers, 28{,}770 (99.0\%) carry a declared or regression-imputed MiFID risk band, and the remaining 320 are excluded from all coherence computations. Customer profiles also include \texttt{customerType} and \texttt{investmentCapacity}, both of which enter the RQ2 regression controls in Section~\ref{sec:setup:stats}.

\subsection{Asset Risk Band Assignment}
\label{sec:setup:bandassignment}

Each asset is assigned to one of the four MiFID risk bands by a three-step hierarchical rule, applied in order.

\subsubsection{Step 1: subcategory metadata} Mutual funds are labelled directly from their subcategory: Money Market funds are treated as Conservative, Bonds as Income, Balanced as Balanced, and Equity or Large Cap as Aggressive. Bonds are labelled by issuer type: Government as Conservative, Corporate as Income, and any remaining bond subcategory as Income by default.

\subsubsection{Step 2: volatility quartiles for stocks} Stocks carry no subcategory tag we can use, so they fall through to a volatility rule. For each ISIN we compute the annualised standard deviation of daily log returns on a trailing 252-day window, then bucket all stocks into the four bands by the stock-only quartile thresholds $(q_1, q_2, q_3)$.

\subsubsection{Step 3: residual default} Any asset that falls through both prior steps is assigned to Balanced.

\paragraph{Resulting distribution.} Applying the three steps above produces an asset-band distribution of $(190, 333, 105, 178)$ across (Conservative, Income, Balanced, Aggressive) over $|A| = 806$ assets. Under the formula defined in Section~\ref{sec:method:pck}, this yields a band-conditional random baseline of $\pi(b) = (0.65, 0.78, 0.76, 0.35)$.

\subsection{Temporal Split Protocol}
\label{sec:setup:splits}

Evaluation uses 69 temporal splits on an evenly spaced trading-day grid between August 2019 and May 2022, yielding approximately biweekly split points. For each split, the training set contains all Buys in $[0, \texttt{time\_point}]$ and the test set contains all Buys in $[\texttt{time\_point}, \texttt{test\_end}]$, deduplicated against training and filtered to the training-asset universe. Eligible customers have at least one interaction in both training and test, while eligible assets have prices on both endpoints. In RQ4 (Section~\ref{sec:results:rq4}), per-split metric deltas are paired against the global LightGCN baseline to quantify the effect of each intervention component.

\subsection{Baselines and Hyperparameter Grids}
\label{sec:setup:baselines}

Table~\ref{tab:baseline-grids} reports the search grids, fixed values, and best-trial selections for the two FAR-Trans baselines. Both grids are swept with Ray Tune, with LightGCN tuned on average nDCG@10 and Random Forest on average ROI@10 across the 69 splits.

\begin{table}[t]
\centering
\small
\begin{tabular}{llll}
\toprule
Parameter & Search space & Best & Role \\
\midrule
\multicolumn{4}{l}{\textbf{LightGCN} (tuned on average nDCG@10)} \\
\texttt{embedding\_dim}    & $\{64, 128\}$          & $64$      & grid \\
\texttt{num\_layers}       & $\{2, 3\}$             & $3$       & grid \\
\texttt{learning\_rate}    & $\{10^{-2}, 10^{-3}\}$ & $10^{-3}$ & grid \\
\texttt{weight\_decay}     & $10^{-5}$              & ---       & fixed \\
\texttt{keep\_prob}        & $0.6$                  & ---       & fixed \\
\texttt{num\_epochs}       & $50$                   & ---       & fixed \\
\texttt{batch\_size}       & $1024$                 & ---       & fixed \\
\midrule
\multicolumn{4}{l}{\textbf{Random Forest} (tuned on average ROI@10)} \\
\texttt{num\_estimators}   & $\{20, 30, 40, 50\}$   & $40$      & grid \\
\texttt{max\_depth}        & $\{15, 25, 50\}$        & $50$      & grid \\
horizon                    & 6 months               & ---       & fixed \\
\texttt{random\_state}     & $42$                   & ---       & fixed \\
\bottomrule
\end{tabular}
\caption{Baseline hyperparameter grids and best-trial values. ``grid'' rows are searched, ``fixed'' rows are held constant.}
\label{tab:baseline-grids}
\end{table}

The stratified PC-LightGCN inherits the best LightGCN trial's backbone hyperparameters. Two trials are evaluated, $\lambda = 0$ (BPR only, isolating the stratified architecture) and $\lambda = 1$ (coherence margin activated), on the same 69 splits.

\subsection{Statistical Inference}
\label{sec:setup:stats}

Three of the four research questions from Section~\ref{sec:intro} require statistical analysis beyond the headline metrics. RQ1 is descriptive and needs no inference test. The remaining three map to the following procedures.

\begin{itemize}
    \item \textbf{RQ2: Does discordance carry a return penalty?} We fit a transaction-level OLS (Equation~\ref{eq:rq2-ols}) on 208{,}029 priced Buys with cluster-robust standard errors on \texttt{customerID} and a two-sided Wald test at $\alpha = 0.05$. Clustering is warranted because each customer contributes roughly 8 Buys on average, so within-customer residuals share unobserved shocks.
    \item \textbf{RQ3: Do baselines lift coherence on every band?} For each (declared band, model) cell, we average per-customer PC@10 unweighted over all (customer, split) pairs. Per-band lift is then the ratio of that cell mean to the random baseline $\pi(b)$, which normalises away the band-specific easiness of the asset universe (Section~\ref{sec:method:pck}).
    \item \textbf{RQ4: Does PC-LightGCN improve coherence, and at what cost?} We average each metric over the 69 splits and compute paired per-split deltas against the baseline LightGCN. The win rate reports the fraction of splits on which the treatment outperforms the comparator.
\end{itemize}

The RQ2 OLS specification is given below and interpreted in Section~\ref{sec:results:rq2}:
\begin{equation}
\label{eq:rq2-ols}
\begin{split}
r ={}& \beta_0 + \beta_1 \cdot \text{is\_coherent} \\
    &{}+ \beta_2 \cdot \text{asset\_volatility} \\
    &{}+ \mathbf{C}(\text{type})\boldsymbol{\gamma} \\
    &{}+ \mathbf{C}(\text{year})\boldsymbol{\delta} \\
    &{}+ \varepsilon.
\end{split}
\end{equation}
where $r = (p_{\text{end}} - p_{\text{start}})/p_{\text{start}}$ is the 6-month realised return as a decimal fraction (e.g., $r = 0.03$ means a 3 percentage-point gain) and \texttt{is\_coherent} is binary ($d \leq 1$). Prices are matched via a backward-nearest join: the start price is the most recent close on or before the trade date, and the end price is the most recent close on or before the date six months later. Transactions without a price match within 7 days at either endpoint are dropped.
